\subsubsection{Unequal intensities}

Suppose $D\neq0$, so $|A|\neq|B|$.  Before solving, use the two phase
changes
\[
 A=\ee^{\ii(\chi+\xi)}\widehat A,
 \qquad B=\ee^{\ii(\chi-\xi)}\widehat B
\]
to replace $\eta_t$ and $\eta_x$ by the nonnegative forcing strengths
\begin{equation}
 \gamma_t=|\eta_t|,
 \qquad \gamma_x=|\eta_x|.
 \label{eq:forcing-strengths}
\end{equation}
The coefficients $a,b,c$ and all intensities are unchanged.  We drop the
hats below.  Thus \eqref{eq:difference-balance} becomes
\begin{equation}
 D(a+bT)-2\ii\gamma_x k=0.
 \label{eq:unequal-difference}
\end{equation}

\paragraph{No input: $\gamma_t=\gamma_x=0$.}
Since $D\neq0$, \eqref{eq:unequal-difference} gives $a+bT=0$.
Substitution into either equilibrium equation then gives $b=c$.
Consequently an unequal-intensity equilibrium occurs only in the degenerate
case
\begin{equation}
 b=c,
 \qquad T=-\frac ab\in\mathbb R_{>0},
 \qquad 0<x<T.
 \label{eq:unforced-unequal}
\end{equation}
Both phases are arbitrary.  For $b\neq c$, there is no unequal-intensity
equilibrium without input.

\paragraph{Only temporal input: $\gamma_t>0$ and $\gamma_x=0$.}
Again $T=-a/b$ must be real and positive.  The remaining equation is
\begin{equation}
 |c-b|^2x(T-x)=\gamma_t^2.
 \label{eq:temporal-unequal-balance}
\end{equation}
For $b\neq c$, the two exchange-related intensities are
\begin{equation}
 x_\pm=
 \frac{T\pm\sqrt{T^2-4\gamma_t^2/|c-b|^2}}{2},
 \qquad |B|^2=T-x_\pm.
 \label{eq:temporal-unequal-strength}
\end{equation}
They exist when
\begin{equation}
 T^2>\frac{4\gamma_t^2}{|c-b|^2}.
 \label{eq:temporal-unequal-existence}
\end{equation}
Writing $A=\sqrt{x}\ee^{\ii\phi}$ and
$B=\sqrt{T-x}\ee^{\ii\psi}$, the temporal input fixes the phase sum by
\begin{equation}
 \ee^{-\ii(\phi+\psi)}
 =-\frac{(c-b)\sqrt{x(T-x)}}{\gamma_t},
 \label{eq:temporal-unequal-phase}
\end{equation}
whereas $\phi-\psi$ remains arbitrary.  This is a biased pair of
counter-propagating waves.

\paragraph{Spatial input, with or without temporal input: $\gamma_x>0$.}
Set $K=k/D\in\mathbb R$.  The real and imaginary parts of
\eqref{eq:unequal-difference} give, when $\RePart b\neq0$,
\begin{equation}
 T=-\frac{\RePart a}{\RePart b},
 \qquad
 K=\frac{\Im(a+bT)}{2\gamma_x}.
 \label{eq:unequal-TK}
\end{equation}
The remaining conditions reduce to
\begin{subequations}
\label{eq:unequal-scalar-system}
\begin{align}
 h^2&=\frac{T^2}{4}-\left(K^2+\frac14\right)D^2,
 \label{eq:unequal-h}\\
 0&=\frac{|c-b|^2}{4}(T^2-D^2)
 +2\gamma_x\RePart(c-b)h
 +(1+4K^2)\gamma_x^2-\gamma_t^2
 +2TK\gamma_x\Im(c-b).
 \label{eq:unequal-balance}
\end{align}
\end{subequations}
Equation \eqref{eq:unequal-balance} is linear in $h$; substituting it into
\eqref{eq:unequal-h} gives a quadratic equation for $D^2$.  Retain only
real roots satisfying
\begin{equation}
 0<D^2<T^2,
 \qquad
 \frac{T^2}{4}-\left(K^2+\frac14\right)D^2\geq0.
 \label{eq:unequal-admissibility}
\end{equation}
For each admissible root, both signs of $D$ give exchange-related states,
and
\begin{equation}
 x=\frac{T+D}{2},
 \qquad k=KD,
 \qquad \zeta=\frac{h-\ii k}{x}.
 \label{eq:unequal-reconstruction}
\end{equation}
If $\gamma_t>0$, the remaining phase is fixed by
\begin{equation}
 u=-\frac{\gamma_x+zeta
 \left[a+(b|\zeta|^2+c)x\right]}{\gamma_t},
 \qquad |u|=1,
 \label{eq:unequal-u}
\end{equation}
and $A=\sqrt{x}\ee^{\ii\phi}$ with $u=\ee^{-2\ii\phi}$ and $B=\zeta A$.
If $\gamma_t=0$, the numerator in \eqref{eq:unequal-u} must vanish and the
common phase $\phi$ is arbitrary.  Thus spatial-only input leaves the
temporal phase free, while combined input generically locks both phases.

The formulas above assume nonzero denominators.  If $\RePart b=0$, solve
the two real equations in \eqref{eq:unequal-difference} directly for
$T$ and $K$.  If $\RePart(c-b)=0$, use
\eqref{eq:unequal-balance} first to determine $D^2$ and then
\eqref{eq:unequal-h} to determine $h$.  A nonpositive $T$ or the absence
of an admissible root means that no unequal-intensity equilibrium exists.

\subsubsection{Equal intensities}

Suppose $D=0$, and write $|A|=|B|=r$.  Equation
\eqref{eq:difference-balance} gives $\eta_xk=0$.

If $\gamma_x=0$, write $A=r\ee^{\ii\phi}$,
$B=r\ee^{\ii\psi}$, and $\Theta=\phi+\psi$.  Then
\begin{equation}
 a+(b+c)r^2+\gamma_t\ee^{-\ii\Theta}=0.
 \label{eq:equal-temporal-balance}
\end{equation}
For $b+c\neq0$, the possible strengths are
\begin{equation}
 r_\pm^2=
 \frac{-\RePart\!\left(a\overline{(b+c)}\right)
 \pm\sqrt{
 \RePart\!\left(a\overline{(b+c)}\right)^2
 -|b+c|^2(|a|^2-\gamma_t^2)}}{|b+c|^2}.
 \label{eq:equal-temporal-strength}
\end{equation}
Only real positive roots are retained.  When $\gamma_t>0$,
\begin{equation}
 \ee^{-\ii\Theta}=-\frac{a+(b+c)r^2}{\gamma_t},
 \label{eq:equal-temporal-phase}
\end{equation}
while $\phi-\psi$ is arbitrary.  When $\gamma_t=0$, one instead requires
$a+(b+c)r^2=0$, and both phases are arbitrary.

If $\gamma_x>0$, then $k=0$ and hence
\begin{equation}
 B=\sigma A,
 \qquad \sigma\in\{1,-1\}.
 \label{eq:equal-parity}
\end{equation}
Set $d=b+c$, $g_\sigma=a+\sigma\gamma_x$, and $y=|A|^2$.  The amplitude
condition is
\begin{equation}
 |g_\sigma+dy|=\gamma_t.
 \label{eq:equal-parity-balance}
\end{equation}
For $d\neq0$, its candidates are
\begin{equation}
 y_{\sigma,\pm}=
 \frac{-\RePart(g_\sigma\overline d)
 \pm\sqrt{\RePart(g_\sigma\overline d)^2
 -|d|^2(|g_\sigma|^2-\gamma_t^2)}}{|d|^2}.
 \label{eq:equal-parity-strength}
\end{equation}
Again, only real positive roots are retained.  If $\gamma_t>0$, the phase
is fixed by
\begin{equation}
 \ee^{-2\ii\phi}=-\frac{g_\sigma+dy}{\sigma\gamma_t}.
 \label{eq:equal-parity-phase}
\end{equation}
If $\gamma_t=0$, require $g_\sigma+dy=0$, and $\phi$ is arbitrary.  The
states $\sigma=1$ and $\sigma=-1$ are the even and odd standing waves,
respectively.

\subsection{One-zero equilibrium solution}

Let $A=0$ and $B=r\ee^{\ii\psi}\neq0$.  The equilibrium equations reduce
to
\begin{equation}
 r^2=-\frac ab\in\mathbb R_{>0},
 \qquad
 \gamma_t\ee^{-2\ii\psi}+\gamma_x=0.
 \label{eq:one-zero-solution}
\end{equation}
Therefore the branch exists with arbitrary $\psi$ when there is no input,
does not exist with only one input, and exists with both inputs precisely
when $\gamma_t=\gamma_x$.  In the last case
$\ee^{-2\ii\psi}=-1$ in the normalized phase coordinates.  The state
$(A,0)$ follows by exchanging $A$ and $B$.  These are left- and
right-traveling waves.

\subsection{Zero equilibrium solution}

The zero equilibrium
\begin{equation}
 (A,B)=(0,0)
 \label{eq:zero-equilibrium}
\end{equation}
always exists.  It has no leading critical-mode activation.

\subsection{Stability}

At an equilibrium $(A_*,B_*)$, the linearization is
\begin{equation}
 \frac{d}{d\tau}
 \begin{pmatrix}\delta\mathbf A\\ \delta\overline{\mathbf A}\end{pmatrix}
 =\mathcal J
 \begin{pmatrix}\delta\mathbf A\\ \delta\overline{\mathbf A}\end{pmatrix},
 \qquad
 \mathcal J=\begin{pmatrix}M&N\\ \overline N&\overline M\end{pmatrix},
 \label{eq:full-Jacobian}
\end{equation}
where $\delta\mathbf A=(\delta A,\delta B)^\mathsf T$ and
\begin{equation}
 M=\begin{pmatrix}
 a+2b|A_*|^2+c|B_*|^2&cA_*\overline{B_*}+\gamma_x\\
 cB_*\overline{A_*}+\gamma_x&a+2b|B_*|^2+c|A_*|^2
 \end{pmatrix},
 \quad
 N=\begin{pmatrix}
 bA_*^2&cA_*B_*+\gamma_t\\
 cA_*B_*+\gamma_t&bB_*^2
 \end{pmatrix}.
 \label{eq:stability-blocks}
\end{equation}
An isolated equilibrium is asymptotically stable if and only if all four
eigenvalues of $\mathcal J$ have negative real part.  Equivalently, if
\begin{equation}
 \det(\lambda I-\mathcal J)
 =\lambda^4+q_1\lambda^3+q_2\lambda^2+q_3\lambda+q_4,
 \label{eq:stability-quartic}
\end{equation}
then
\begin{equation}
 q_j>0\quad(j=1,\ldots,4),
 \qquad q_1q_2>q_3,
 \qquad q_1q_2q_3>q_3^2+q_1^2q_4.
 \label{eq:stability-Hurwitz}
\end{equation}
Substitution of \eqref{eq:unequal-reconstruction} or
\eqref{eq:one-zero-solution} into \eqref{eq:stability-blocks} therefore
gives the stable unequal-intensity and one-zero equilibria without further
case-dependent variables.

For an equal-intensity state $B_*=\sigma A_*$, perturbations split into
$\delta B=\kappa\delta A$, $\kappa\in\{1,-1\}$.  Define
\begin{align}
 P_{\sigma\kappa}
 &=a+\kappa\gamma_x
 +(2b+c+\kappa\sigma c)|A_*|^2,
 \label{eq:equal-P}\\
 Q_{\sigma\kappa}
 &=(b+\kappa\sigma c)A_*^2+\kappa\gamma_t.
 \label{eq:equal-Q}
\end{align}
The two eigenvalues in this block satisfy
\begin{equation}
 \lambda^2-2\RePart(P_{\sigma\kappa})\lambda
 +|P_{\sigma\kappa}|^2-|Q_{\sigma\kappa}|^2=0.
 \label{eq:equal-characteristic}
\end{equation}
With both inputs, the equilibrium is asymptotically stable precisely when,
for $\kappa=\pm1$,
\begin{equation}
 \RePart(P_{\sigma\kappa})<0,
 \qquad
 |P_{\sigma\kappa}|^2>|Q_{\sigma\kappa}|^2.
 \label{eq:equal-stability}
\end{equation}
With spatial input only, the $\kappa=\sigma$ block contains the neutral
common-phase mode.  Orbital stability requires
\begin{equation}
 |P_{\sigma\sigma}|^2=|Q_{\sigma\sigma}|^2,
 \qquad \RePart(P_{\sigma\sigma})<0,
 \label{eq:spatial-orbital}
\end{equation}
together with \eqref{eq:equal-stability} for $\kappa=-\sigma$.  With
temporal input only, use the free spatial phase to take $B_*=A_*$; the
$\kappa=-1$ block is neutral, and the analogous orbital conditions apply.

Without input, the standing wave $|A_*|=|B_*|=r$ is orbitally stable if and
only if
\begin{equation}
 \RePart(b+c)<0,
 \qquad \RePart(b-c)<0.
 \label{eq:unforced-standing-stability}
\end{equation}
The one-zero traveling wave, with $r^2=-a/b>0$, is orbitally stable if and
only if
\begin{equation}
 \RePart b<0,
 \qquad \RePart(a+cr^2)<0.
 \label{eq:unforced-traveling-stability}
\end{equation}

At the zero equilibrium, $Z_\sigma=A+\sigma B$ satisfies
\begin{equation}
 (Z_\sigma)_\tau=(a+\sigma\gamma_x)Z_\sigma
 +\sigma\gamma_t\overline{Z_\sigma}.
 \label{eq:zero-parity}
\end{equation}
Hence the zero equilibrium is stable if and only if
\begin{equation}
 \RePart(a+\sigma\gamma_x)<0,
 \qquad
 |a+\sigma\gamma_x|^2>\gamma_t^2,
 \qquad \sigma=\pm1.
 \label{eq:zero-stability}
\end{equation}
For no input this reduces to $\RePart a<0$.

\subsection{Input--pattern summary}

From \eqref{eq:u1}, a one-zero state is a traveling wave, an
equal-intensity state is a standing wave, and an unequal-intensity state is
a biased counter-propagating wave.  The leading complex-mode strengths are
summarized below; the physical field also contains the complex conjugate.
\begin{table}[ht]
\centering
\small
\begin{tabular}{@{}llll@{}}
\toprule
Input & Existence condition & Activation pattern & Strength \\
\midrule
None
& $-a/b=r^2>0$
& traveling wave
& one component $r$ \\
None
& $-a/(b+c)=r^2>0$
& standing wave
& envelope $2r$ \\
Temporal only
& \eqref{eq:equal-temporal-strength}
& time-locked standing wave
& envelope $2r_\pm$ \\
Temporal only
& \eqref{eq:temporal-unequal-strength}
& biased counter-propagating wave
& $\sqrt{x_\pm},\sqrt{T-x_\pm}$ \\
Spatial only
& $-(a+\sigma\gamma_x)/(b+c)=r^2>0$
& even/odd pinned standing wave
& envelope $2r$ \\
Both
& \eqref{eq:equal-parity-strength}
& fully locked standing wave
& envelope $2\sqrt{y_{\sigma,\pm}}$ \\
Both, $\gamma_t=\gamma_x$
& $-a/b=r^2>0$
& locked traveling wave
& one component $r$ \\
Spatial/both
& \eqref{eq:unequal-scalar-system}
& biased counter-propagating wave
& $|A|,|B|$ \\
Any
& \eqref{eq:zero-stability}
& stable inactive state
& $0$ \\
\bottomrule
\end{tabular}
\caption{Input-induced activation patterns.  An existing nonzero branch is
stable only when it also satisfies \eqref{eq:stability-Hurwitz},
\eqref{eq:equal-stability}, \eqref{eq:unforced-standing-stability}, or
\eqref{eq:unforced-traveling-stability}, as appropriate.}
\label{tab:input-patterns}
\end{table}
